\documentclass[10pt,a4paper]{amsart}
\usepackage[margin=19mm]{geometry}
\usepackage{amsmath,amssymb,mathtools}
\usepackage[scr=rsfs,cal=euler]{mathalfa}
\usepackage{xfrac}
\usepackage[unicode,hidelinks,bookmarks=false]{hyperref}
\newcommand{\ie}{i.e.\ }
\newcommand{\cf}{cf.\ }
\newcommand{\resp}{respectively\ }
\newcommand{\Subsection}{\subsection}
\providecommand{\nobreakdash}{\nobreak\hbox{-}}
\setlength{\emergencystretch}{2em}
\multlinegap0pt
\usepackage{bm}
\usepackage{bbm}
\usepackage{dsfont}
\usepackage{xspace}
\usepackage{cancel}
\usepackage{mathtools}
\usepackage{amsthm}
\makeatletter
\@ifundefined{theorem}{\newtheorem{theorem}{Theorem}}{}
\@ifundefined{lemma}{\newtheorem{lemma}[theorem]{Lemma}}{}
\makeatother
\newcommand{\Mod}[1]{\ (\mathrm{mod}\ #1)}
\title[Partitions into Sums of Two Squares]{Partitions into Sums of Two Squares}
\author{Jaime Palacios}
\date{Source: arXiv:2508.17662v1 (CC BY 4.0)}
\begin{document}
\maketitle

\noindent\emph{Source and licence.} Partitions into Sums of Two Squares. Version arXiv:2508.17662v1 (CC BY 4.0), distributed under the Creative Commons Attribution 4.0 International licence, as recorded in the arXiv licence metadata for this exact version and shown on the abstract page https://arxiv.org/abs/2508.17662v1. Full rights notice: licenses/arxiv-2508.17662-CC-BY-4.0.txt.\par\smallskip

\noindent\textit{Editorial scope.} Counted unit: the Lemma whose source label is \texttt{deriv} in the article (source0.tex lines 275--280, source label \texttt{deriv}), delivered together with the article's own complete original proof of it (source0.tex lines 282--343, one \texttt{proof} environment of 62 lines). No mathematics is written, repaired or re-derived: statement and proof are the article's own text line for line. Required context retained uncounted: tst (lines 500--502); H (lines 524--530). Every definition, display and label of those lines is retained unchanged. Omitted from this package and not redistributed: 1--274, 281--281, 344--499, 503--523, 531--593. Nothing inside a retained excerpt is omitted or altered: every retained body is delivered exactly as the article's lines are written, so no omission receipt of any kind accompanies this package. Citations: the retained excerpts carry no citation command at all, so no bibliography entry of the article is transcribed and no reference is redistributed. Reference adaptation: none was needed and none was made; every reference inside the delivered unit resolves inside it, so no unresolved reference or citation is produced. Printed slips kept literal: no printed word, symbol, hypothesis or number of the counted statement or of its proof was altered. Layout and class: the article's own class is \texttt{amsart} with packages including inputenc, amssymb, amsfonts, amsmath, amsthm, mathtools, indentfirst, bigints, geometry, xcolor, nccmath, fullpage, fancyhdr, stmaryrd; the wrapper is the lane's pinned \texttt{amsart} editorial wrapper at 10pt on A4, which restates the theorem-like environments the retained text needs and adds the article's own macro definitions that the retained text needs (\texttt{\textbackslash Mod}), with extra wrapper packages (bm, bbm, dsfont, xspace, cancel, mathtools). The delivered numbering is the unit's own. Assets: no retained span contains a figure environment, a graphics-inclusion command, a PDF-page inclusion, a table or a TikZ picture; no image file is copied into this package, the package declares no asset dependency and no image file is redistributed with it. Licence: version arXiv:2508.17662v1 of this article is distributed under the Creative Commons Attribution 4.0 International licence, as recorded in the arXiv licence metadata for this exact version and shown on the abstract page https://arxiv.org/abs/2508.17662v1. A full rights notice is delivered at licenses/arxiv-2508.17662-CC-BY-4.0.txt. Characters: every retained excerpt is pure ASCII, so the delivered engine, XeLaTeX with its default Latin Modern text font, reports no missing character.
\begin{theorem}[Two Square Theorem] \label{tst}
    A positive integer $n$ is the sum of two squares if and only if each of its prime factors $p$ such that $p\equiv 3 \ (\text{mod } 4)$ occurs an even number of times.
\end{theorem}

\begin{theorem}[Hankel] \label{H}
    For any complex number $s$,
    \begin{equation*}
        \frac{1}{2\pi i} \int\limits_\mathcal{H} e^zz^{-s} dz= \frac{1}{\Gamma(s)},
    \end{equation*}
    where for $r>0$ we let $\mathcal{H}=\mathcal{H}(r)$ denote the Hankel contour, which consists  of a path that passes from $-\infty-ir$ to $-ir$ along a straight line, and then from $-ir$ to $ir$ along the semicircle $re^{i\theta}$, $-\frac{\pi}{2} \leq \theta \leq \frac{\pi}{2}$, and then from $ir$ to $- \infty +ir $ along a straight line . Here $z^{-s}$ is assumed to have its principal value.
\end{theorem}

\begin{lemma}\label{deriv} Let $\rho=e^{-1/X}$, then
    \begin{equation} 
        \Big(\rho \frac{d}{d\rho}\Big)^m \Phi(\rho)= \frac{K \zeta(2) \Gamma(m+1)X^{m+1}}{\sqrt{\log X}}\Bigg(1+O\Big(\frac{1}{\log X}\Big)\Bigg),
    \end{equation}
    as $\rho \to 1^-$.
\end{lemma}

\begin{proof}
We have that
\begin{align*}
    \Big(\rho \frac{d}{d\rho}\Big)^m\Phi(\rho)&=\sum_{j=1}^\infty \sum_{\ell\in \mathcal{S}} \frac{1}{j} (j\ell)^m \rho^{j\ell}\\
    & =\sum_{j=1}^\infty \sum_{\ell\in \mathcal{S}} \frac{1}{j} (j\ell)^m e^{-j\ell/X}.
\end{align*}
   Using an inverse Mellin transform and interchanging the sums with the integrals, we obtain
\begin{equation*}
     \Big(\rho \frac{d}{d\rho}\Big)^m\Phi(\rho) = \frac{1}{2\pi i} \int\limits_{c_0-i\infty}^{c_0+i\infty} X^s\zeta(s-m+1)\Gamma(s)\Bigg(\sum_{\ell \in \mathcal{S}} \frac{1}{\ell^{s-m}}\Bigg) ds,
\end{equation*}
for some $c_0>m+1$.

      Theorem \ref{tst} implies that 
    \begin{equation*}
        \sum_{\ell\in \mathcal{S}} \frac{1}{\ell^s}= f(s)\sqrt{\zeta(s)L(s,\chi)},
    \end{equation*}
    where $\chi$ is the non-principal Dirichlet character modulo 4, and 
    \begin{equation*}
        f(s)=(1-2^{-s})^{-1/2}\prod_{p \equiv 3 \Mod 4} (1-p^{-2s})^{-1/2},
    \end{equation*}
which converges for $\Re( s)>\frac{1}{2}$.

Defining the  function $$g(s):=f(s)\sqrt{(s-1)\zeta(s)L(s,\chi_1)},$$  we can then write for some $c_0>m+1$ 
\begin{align*}
    \Big(\rho \frac{d}{d\rho}\Big)^m\Phi(\rho) & =  \frac{1}{2\pi i} \int\limits_{c_0-iT}^{c_0+iT} X^s\zeta(s-m+1)\Gamma(s) \frac{g(s-m)}{\sqrt{s-m-1}} ds+O\Big(\frac{X^{m+1+\varepsilon}}{e^T}\Big)  \quad  \\ %\quad(T>X)
    & =  \frac{1}{2\pi i} \int\limits_{\mathcal{C}} X^s\zeta(s-m+1)\Gamma(s) \frac{g(s-m)}{\sqrt{s-m-1}} ds+O\Big(X^{m+1-c}\log T\Big)  \\
    & =  \frac{1}{2\pi i} \int\limits_{\mathcal{C}} X^s\zeta(s-m+1)\Gamma(s) \frac{g(s-m)}{\sqrt{s-m-1}} ds+O\Bigg(\frac{X^{m+1}}{(\log X)^{3/2}}\Bigg), \\
\end{align*}
where $\mathcal{C}$ is the contour running from $m+1-c-i\delta$ along a straight line to $m+1-i\delta$, then along the semicircle $m+1+\delta e^{i\theta}$, $-\frac{\pi}{2}\leq \theta \leq \frac{\pi}{2}$, and finally along a straight line to $m+1-c+i\delta$. Here we have taken  $c= \frac{c'}{\log T}$ for some constant $c'>0$ so that the contour $\mathcal{C}$ is inside the zero free region for $\zeta(s-m)$ up to height $T=\exp(\sqrt{\log X})$, and we take $\delta= \frac{1}{\log X}$. 

Since $g(s-m)$ has a removable fractional singularity in a neighbourhood of $m+1$, we see that

\begin{align*}
         \frac{1}{2\pi i} \int\limits_{\mathcal{C}} & X^s\zeta(s-m+1)\Gamma(s) \frac{g(s-m)}{\sqrt{s-m-1}} ds \\
         &=  \frac{1}{2\pi i} \int\limits_{\mathcal{C}} X^s\zeta(s-m+1)\Gamma(s) \frac{1}{\sqrt{s-m-1}}\Big(g(1)+O(|s-(m+1)|)\Big) ds\\
    &=  \frac{1}{2\pi i} \int\limits_{\mathcal{C}} X^s\zeta(s-m+1)\Gamma(s) \frac{1}{\sqrt{s-m-1}}\Big(K\sqrt{\pi}+O(|s-(m+1)|)\Big) ds.
\end{align*}
The choice of $\delta=\frac{1}{\log X}$ implies that 

\begin{align*}
   \int\limits_{\mathcal{C}} X^s\zeta(s-m+1)\Gamma(s) \frac{1}{\sqrt{s-m-1}}|s-(m+1)| ds& \ll \int\limits_{\mathcal{C}} X^{\Re s} |s-(m+1)|^{\frac{1}{2}} |ds|\\
   & \ll \frac{X^{m+1}}{(\log X)^{3/2}}.
\end{align*}

Now we note that after a change of variable,
\begin{align*}
    \frac{1}{2\pi i} \int\limits_{\mathcal{C}} X^s\zeta(s-m+1)\Gamma(s) \frac{1}{\sqrt{s-m-1}} ds &= \frac{1}{2\pi i} \int\limits_{\mathcal{C'}} X^{m+1+u}\zeta(2+u)\Gamma(m+1+u) \frac{du}{\sqrt{u}} \\
    & =  \frac{1}{2\pi i} \int\limits_{\mathcal{C'}} X^{m+1+u}\Big(\zeta(2)\Gamma(m+1)+O(u) \Big) \frac{du}{\sqrt{u}},
\end{align*}
where $\mathcal{C'}$ is the contour running from $-c-i\delta$ along a straight line to $-i\delta$, then along the semicircle $\delta e^{i\theta}$, $-\frac{\pi}{2}\leq \theta \leq \frac{\pi}{2}$, and finally along a straight line to $-c+i\delta$. Here we have that 
\begin{align*}
    \frac{1}{2\pi i} \int\limits_{\mathcal{C'}}X^u|u|^{1/2}du&\ll \frac{1}{\sqrt{\log X}} \int\limits_{\mathcal{C'}} |du| \ll \frac{1}{(\log X)^{3/2}},
\end{align*}
and by Theorem \ref{H},
\begin{align*}
    \frac{1}{2\pi i} \int\limits_{\mathcal{C'}}X^u\frac{du}{\sqrt{u}} &=\frac{1}{2\pi i} \int\limits_{\mathcal{H}}X^u \frac{du}{\sqrt{u}}  +O\Big(X^{-c}\Big)\\
    & = \frac{1}{\sqrt{\pi \log X}} +O\Big(\frac{1}{(\log X)^{3/2}}\Big).
\end{align*}
Putting all above together the lemma follows.


\end{proof}

\end{document}
